\section{Causica: a generative reference with explicit admission boundaries}
\label{exec14:causica-contract}

An analyst wants to simulate a larger servicing effort for a borrower, compare
it with the observed treatment, and then aggregate the result into a contact
policy. A learned structural model makes these operations convenient: one
model can represent a graph, conditional mechanisms, noise and interventions.
Convenience does not settle whether the graph is identified, whether a
counterfactual refers to the same borrower, or whether the reported uncertainty
includes the relevant source of error. Causica is useful precisely because
its explicit abstractions let the proposal examine these questions together.

The selected source is Microsoft's Causica revision \texttt{c3b7dd2b09fe},
dated 24 August 2026, declaring package version 0.4.5. Its DECI path combines
an approximate distribution over graphs with learned functional relationships
and noise modules. This is a proposed reference for a structural-model
component. It is not an enabled estimator in the local six-operation service.
The distinction matters when comparing the proposal with an existing library:
describing a model in a monograph does not implement its training, inference,
diagnostics or failure handling. \citep{exec14_causica_primary,exec14_causica_decision}

\subsection{What the original result does and does not justify}

Geffner et al. formulate DECI using a nonlinear additive-noise model with no
latent confounders. Its graph distribution and function/noise parameters are
learned by maximizing an evidence lower bound. Their Theorem 1 is an
infinite-data result for a maximizing solution. Appendix A assumes minimality,
structural identifiability, correct specification, causal sufficiency and
integrability of the log likelihood. Structural identifiability means that
another graph in the admitted model family cannot generate the same observed
distribution. It is an assumption with substantive content, not a property
created by calling the optimizer. \citep{exec14_causica_geffner}

These conditions expose four separate obligations for a banking use. The
variable set must represent the relevant common causes; the structural family
must distinguish the intervention query from observational alternatives;
the fitted mechanisms must approximate the admitted law; and the numerical
procedure must achieve an adequate solution. Universal function approximation
does not verify the first two obligations or make a finite training run a
global maximum. The paper explicitly says that the theorem does not cover
its mixed-type and missing-value extensions. A practical implementation may
support those inputs without inheriting that theorem for them.
\citep{exec14_causica_geffner}

After fitting, graph-averaged intervention simulation answers a model-based
question. In the selected source, graph draws share the learned functional
and noise modules. Those draws are therefore not samples from an independently
specified posterior over all function parameters. Their variation is also
not a frequentist confidence region for the true graph. The product should
report the fitted model, graph distribution and numerical simulation error
separately, keeping any unresolved structural ambiguity visible.
\citep{exec14_causica_primary}

\subsection{Conditioning, intervention and the same borrower's counterfactual}

The pinned \texttt{SEM} interface distinguishes \texttt{do},
\texttt{condition}, noise generation and noise reconstruction. Its selected
conditioning method supports a set only when that set has no parents outside
it; otherwise it raises an unsupported-operation error. This restriction is
meaningful. The paper's more general CATE discussion uses regression on
simulated interventional samples and assumes that treatment does not cause
the conditioning variables. That description should not be silently mapped
to an unrestricted call to the selected conditioning helper.
\citep{exec14_causica_primary,exec14_causica_geffner}

For an exact illustration, let \(X=U_X\) be a standardized pre-action risk
factor, \(T=X+U_T\) a continuous servicing intensity, and
\[
       Y=20T+10X+U_Y
\]
the net recovery before the separately considered cost of changing intensity.
Suppose the three noises are independent centered Gaussians, with variances
\(1,1,100\), and measure \(Y\) in dollars. This is an authored structural
model, not a fitted bank law or an example satisfying general graph
identifiability. Since \(E[X\mid T=t]=t/2\),
\[
 E[Y\mid T=t]=25t,\qquad E[Y\mid\operatorname{do}(T=t)]=20t.
\]
Replacing an unsupported conditioning request by an intervention changes the
answer. Reversing that substitution changes it as well. The first request
selects borrowers through their observed intensity and risk; the second
changes intensity while retaining the population distribution of risk.
\citep{exec14_causica_exact}

Consider a borrower whose factual measurements are \((X,T,Y)=(1,2,55)\).
Reconstructing the noises gives \((U_X,U_T,U_Y)=(1,1,5)\). Holding these
noises fixed gives counterfactual
recoveries of USD~15 at intensity zero and USD~35 at intensity one. The
individual contrast is USD~20. Replacing factual noise by fresh population
noise gives marginal means zero and twenty instead. The contrast happens
to agree in this constant-effect construction, while the borrower's two
counterfactual levels do not. Agreement of an effect difference is therefore
an insufficient test of factual-unit preservation. The pinned \texttt{ite}
helper reconstructs noise once and reuses it for the two interventions;
the \texttt{counterfactual} helper follows the same abduction/action/prediction
logic. \citep{exec14_causica_primary,exec14_causica_exact}

The ATE helper takes a different route: it samples each intervention arm and
subtracts component means. With independent samples of size \(n\) from each
arm in the stipulated model, the Monte Carlo variance of that contrast is
\(400/n\), because each arm has residual variance \(200\). At \(n=100\)
the variance is four square dollars. Coupling the two arms through the same
noise gives an exactly constant twenty-dollar contrast here. Such coupling
can reduce simulation variance, but its appropriateness must be checked for
the intended estimator. Neither calculation establishes sampling uncertainty
for learned functions, uncertainty about the graph, or causal truth.
\citep{exec14_causica_exact}

\subsection{A training penalty is not an output DAG check}

The selected DECI module applies a matrix-exponential acyclicity penalty during
training. Its ENCO adjacency distribution nevertheless samples independent
Bernoulli entries off the diagonal; the selected SEM-distribution path wraps
these matrices without a separate DAG filter. At the initialized zero logits,
each directed edge has probability \(\tfrac12\tfrac12=\tfrac14\).
For two nodes, both opposite edges are present with probability
\(1/16\). This is an exact statement about the initialized distribution,
not an observed failure frequency after training. A penalty may greatly
reduce cyclic mass; its existence alone does not certify every returned
graph. \citep{exec14_causica_primary,exec14_causica_exact}

The numerical consequence is visible without running a neural model.
The selected structural simulator starts at zero and performs one simultaneous
update per node. For a DAG, induction over parent depth explains why a finite
number of such updates can propagate a fixed noise vector to every node.
That argument depends on acyclicity. Consider instead the authored equations
\[
        x_1=1+\tfrac12x_2,\qquad x_2=\tfrac12x_1.
\]
Two updates from zero produce \((1,0)\) and then \((1,1/2)\).
The latter has equation residual \((-1/4,0)\); the algebraic fixed point is
\((4/3,2/3)\). Thus a finite update result can look numerically ordinary while
failing even the supplied cyclic equations. The fixed point is a numerical
comparison, not an independently justified equilibrium causal interpretation.
The proposal should refuse an unqualified cyclic draw rather than promote
its finite update into a valid intervention result.
\citep{exec14_causica_primary,exec14_causica_exact}

Hard knowledge constraints need the same admission check. The pinned source
uses zero for prohibited edges, one for required edges and NaN for unspecified
edges. Shape checks and diagonal masking do not by themselves prevent a
required directed cycle. A public contract should use explicit typed states,
validate their values and causal consistency, and translate them into this
internal representation. Adjacency entries also need a declared
parent-to-child convention; they must not be confused with signed structural
coefficients. If an adapter discards cyclic samples, it changes the effective
graph distribution by conditioning on acceptance. It must retain the rejected
mass and describe any reweighting, so a failure cannot disappear from the
denominator. \citep{exec14_causica_primary,exec14_causica_decision}

\subsection{A vector score needs a vector reference}

The evaluation boundary has another concrete issue. The pinned ATE helper
averages over observations and retains the outcome's component dimension.
The selected \texttt{eval\_ate\_rmse} reference instead applies
\texttt{.mean()} without specifying an axis to each intervention dataset.
For a multidimensional node this collapses observations and components
together before subtracting the two arms. The resulting scalar is then
compared with a component vector. The source can therefore assign a nonzero
error to an exact vector estimate. \citep{exec14_causica_primary}

Take two equal-size reference arms whose component means are \((2,8)\) and
\((0,0)\), with both components measured in dollars. The correct vector
effect is \((2,8)\). An exact prediction has zero componentwise error, but
the all-axis reference is \(5\), producing
\[
       \sqrt{(2-5)^2+(8-5)^2}=\sqrt{18}.
\]
This is an independent exact trace of the inspected reduction, not the output
of an executed Causica test. The selected integration test checks result shape;
that assertion alone does not detect this counterexample. The pinned pytest
configuration also excludes the integration directory from default recursive
collection, so the existence of a CI workflow running the general test tree
is not evidence that this integration test ran.
\citep{exec14_causica_primary,exec14_causica_exact}

The repair route is specific. First preserve the observation axis and calculate
each component's treatment contrast. Then declare whether the requested score
is a vector, a weighted monetary sum or a norm. Different currencies, horizons,
one-hot categories or scales must not be pooled implicitly. An executable
prototype should supply a deterministic SEM or fixed prediction reference,
check zero error for an exact multicomponent effect, and test deliberate
perturbations component by component. Only after these checks should noisy
simulation be added. This separates a scorer error from simulation variation
and model error, reducing the chance of rejecting a correct estimator or
rewarding an incorrect one. \citep{exec14_causica_decision}

\subsection{Graph averaging must preserve the decision risk}

A useful numerical average may still conceal a structurally fragile decision.
Suppose two admitted models imply effects of USD~20 and minus USD~20 and have
stipulated weights \(3/4\) and \(1/4\). The weighted effect is USD~10 and the
between-model variance is \(300\) square dollars. Averaging \(R\) independent
model draws estimates that weighted mean with Monte Carlo variance \(300/R\).
As \(R\) grows, numerical error vanishes around the chosen weighted mean;
the two component effects remain different. The weights themselves are not
validated by increasing the simulation count. \citep{exec14_causica_exact}

With an assumed incremental cost of USD~8, the component net values are
USD~12 and minus USD~28, while their weighted average is USD~2. A manager who
requires robustness across the admitted models cannot infer it from that
positive weighted average. The output should show the component consequences
and the basis for the weights, alongside a model-averaged recommendation
when one is requested. If observationally compatible models remain relevant,
the next information-gathering decision concerns which restriction or
experiment could distinguish their action consequences. A narrow numerical
error bar around an average does not answer that question.
\citep{exec14_causica_exact,exec14_causica_decision}

\subsection{Reuse choice and a bounded prototype}

The pinned root license is MIT. The package declares Python~3.10 and
Torch~2.0.0; its saved lock contains 277 packages, including Lightning,
TensorDict, NumPy and MLflow. These records provide a reproducible dependency
starting point, not proof that an installation or every dependency right has
been qualified. The saved repository metadata does not mark it archived,
but neither an activity timestamp nor a CI badge is a support commitment.
The README directs users of older Rhino and DDECI paths to a previous release.
The current tree contains adjacent temporal classes, yet the selected
intervention helper explicitly refuses temporal functional relationships.
These are distinct paths requiring distinct reviews.
\citep{exec14_causica_primary}

Direct adoption is deferred because it would inherit unresolved graph,
conditioning, scoring and environment boundaries. A wholesale port is not
selected: it would transfer neural, noise-distribution, optimization and
dependency maintenance without demonstrated benefit. The selected near-term
use is a pinned reference, followed by a possible narrow adapter that preserves
upstream mechanics and validates the contract at its boundary. A supplied
graph simulator and a graph-class effect analysis remain substantive
comparators; they expose different tradeoffs between structural commitments,
flexibility and uncertainty. \citep{exec14_causica_decision}

The first prototype can use public or generated mechanisms. It should check
known DAG propagation, reversed node order, intervention versus conditioning,
factual-noise preservation, scalar and vector effects, required cycles, failed
optimization, repeated interventions, temporal refusals and batch dimensions.
It should retain unsuccessful requests and separate fitted-model error,
simulation error and structural ambiguity. Training and simulation time,
dependency work, expert review, conversion failures and update maintenance
belong in the cost comparison. Bank experiments become justified only after
this bounded work demonstrates that the proposed integration preserves a
useful causal question and improves a declared workflow. No trained-model
performance, bank return or enabled integration is inferred from the source
review or exact examples. \citep{exec14_causica_decision}

\label{exec14:causica-contract-end}
